% ═══════════════════════════════════════════════════════════════
% ARBEITSBLATT: Elektromagnetische Induktion (Kl. 9, DS 09b)
% Begleitet LP-09b
% ═══════════════════════════════════════════════════════════════

\documentclass[11pt, a4paper]{article}

\usepackage[utf8]{inputenc}
\usepackage[T1]{fontenc}
\usepackage[ngerman]{babel}

% --- SCHRIFTEN ---
\usepackage{helvet}
\renewcommand{\familydefault}{\sfdefault}

\usepackage{microtype}
\usepackage[margin=1.5cm, top=1.8cm, bottom=1.5cm]{geometry}
\usepackage{enumitem}
\usepackage{tabularx}
\usepackage{booktabs}
\usepackage{array}
\usepackage{tikz}

\usepackage{amsmath, amssymb}
\usepackage{siunitx}
\sisetup{locale = DE, per-mode = symbol}

\usepackage{fancyhdr}
\usepackage{lastpage}

\usepackage{xcolor}
\usepackage{tcolorbox}
\tcbuselibrary{skins}

% ═══════════════════════════════════════════════════════════════
% FARBEN
% ═══════════════════════════════════════════════════════════════

\definecolor{physik}{HTML}{2c5aa0}
\definecolor{hellgrau}{HTML}{F5F5F5}
\definecolor{mittelgrau}{HTML}{888888}

\colorlet{fachfarbe}{physik}

% ═══════════════════════════════════════════════════════════════
% VARIABLEN
% ═══════════════════════════════════════════════════════════════

\newcommand{\fach}{Physik}
\newcommand{\thema}{Elektromagnetische Induktion}
\newcommand{\klasse}{9}

% ═══════════════════════════════════════════════════════════════
% KOPF- UND FUSSZEILE
% ═══════════════════════════════════════════════════════════════

\pagestyle{fancy}
\fancyhf{}
\renewcommand{\headrulewidth}{0.4pt}
\renewcommand{\footrulewidth}{0pt}
\lhead{\small \fach{} Klasse \klasse{} | \thema}
\rhead{\small Name: \underline{\hspace{3.5cm}}}
\cfoot{\small\textcolor{mittelgrau}{Seite \thepage{} von \pageref{LastPage}}}

% ═══════════════════════════════════════════════════════════════
% BOXEN
% ═══════════════════════════════════════════════════════════════

\newtcolorbox{merkkasten}[1][]{
    colback=white,
    colframe=fachfarbe,
    coltitle=black,
    colbacktitle=white,
    boxrule=1pt,
    arc=0pt,
    left=8pt, right=8pt, top=6pt, bottom=5pt,
    before skip=5pt, after skip=5pt,
    fonttitle=\bfseries,
    title=#1
}

\newtcolorbox{wortbank}{
    colback=white,
    colframe=fachfarbe,
    boxrule=1pt,
    arc=0pt,
    left=6pt, right=6pt, top=4pt, bottom=4pt,
    before skip=3pt, after skip=3pt
}

% ═══════════════════════════════════════════════════════════════
% BEFEHLE
% ═══════════════════════════════════════════════════════════════

\newcommand{\luecke}[1]{\underline{\hspace{#1}}}

\setlength{\parindent}{0pt}
\setlength{\parskip}{3pt}

% ═══════════════════════════════════════════════════════════════
% DOKUMENT
% ═══════════════════════════════════════════════════════════════

\begin{document}

% ═══════════════════════════════════════════════════════════════
% HEADER
% ═══════════════════════════════════════════════════════════════

\begin{center}
    {\Large\bfseries\textcolor{fachfarbe}{\thema}}\\[0.15em]
    {\small Arbeitsblatt | \fach{} Klasse \klasse}
\end{center}
\vspace{-0.4em}
\hrule
\vspace{0.5em}

% ═══════════════════════════════════════════════════════════════
% KASTEN 1: WAS IST INDUKTION?
% ═══════════════════════════════════════════════════════════════

\begin{merkkasten}[K1: Was ist Induktion?]

\begin{wortbank}
\centering
\textbf{Bewegung} ~ | ~ \textbf{Spannung} ~ | ~ \textbf{Spule} ~ | ~ \textbf{Magneten} ~ | ~ \textbf{keine} ~ | ~ \textbf{Faraday}
\end{wortbank}

\vspace{2mm}

Bewegt man einen \luecke{2.5cm} durch eine \luecke{2.5cm}, so entsteht eine elektrische \luecke{2.5cm}.

\vspace{2mm}

Diesen Vorgang nennt man \textbf{elektromagnetische Induktion}.

\vspace{2mm}

\textbf{Wichtig:} Induktion braucht \luecke{2.5cm} (Änderung). Ohne Bewegung gibt es \luecke{2.5cm} Spannung!

\vspace{2mm}

Entdeckt wurde die Induktion 1831 von Michael \luecke{2.5cm}.

\end{merkkasten}

% ═══════════════════════════════════════════════════════════════
% KASTEN 2: EINFLUSSFAKTOREN
% ═══════════════════════════════════════════════════════════════

\begin{merkkasten}[K2: Wovon hängt die Spannung ab?]

\textbf{Beobachte in der Simulation und ergänze:}

\vspace{2mm}

\begin{tabularx}{\textwidth}{@{}X|l@{}}
Wenn man den Magneten \textbf{schneller} bewegt... & ...wird die Spannung \luecke{2.5cm}. \\[0.4em]
\hline
Wenn man einen \textbf{stärkeren} Magneten verwendet... & ...wird die Spannung \luecke{2.5cm}. \\[0.4em]
\hline
Wenn die Spule \textbf{mehr Windungen} hat... & ...wird die Spannung \luecke{2.5cm}. \\[0.4em]
\end{tabularx}

\vspace{3mm}

\textbf{Was passiert, wenn man die Bewegungsrichtung umkehrt?}

\vspace{1mm}
Die \luecke{2.5cm} der Spannung ändert sich (das \luecke{2.5cm} wechselt).

\end{merkkasten}

% ═══════════════════════════════════════════════════════════════
% KASTEN 3: MOTOR UND GENERATOR
% ═══════════════════════════════════════════════════════════════

\begin{merkkasten}[K3: Motor und Generator -- Geschwister]

\begin{wortbank}
\centering
\textbf{Bewegungsenergie} ~ | ~ \textbf{Elektrische Energie} ~ | ~ \textbf{hinein} ~ | ~ \textbf{heraus}
\end{wortbank}

\vspace{2mm}

Motor und Generator sind wie \textbf{Geschwister} -- sie haben den gleichen Aufbau, aber funktionieren \textbf{umgekehrt}.

\vspace{3mm}

\begin{tabularx}{\textwidth}{@{}l|X|X@{}}
& \textbf{Motor} & \textbf{Generator} \\
\hline
Strom fließt... & \luecke{2cm} & \luecke{2cm} \\[0.3em]
\hline
Wandelt um: & \luecke{3cm} $\rightarrow$ Bewegung & Bewegung $\rightarrow$ \luecke{3cm} \\[0.3em]
\end{tabularx}

\vspace{3mm}

\textbf{Beispiele für Generatoren:} \luecke{6cm}

\end{merkkasten}

% ═══════════════════════════════════════════════════════════════
% ÜBUNG 1: ANWENDUNGEN
% ═══════════════════════════════════════════════════════════════

\begin{merkkasten}[Ü1: Anwendungen zuordnen]

\textbf{Kreuze an: Motor oder Generator?}

\vspace{2mm}

\begin{tabularx}{\textwidth}{@{}X|c|c@{}}
\textbf{Gerät} & \textbf{Motor} & \textbf{Generator} \\
\hline
Fahrraddynamo & $\square$ & $\square$ \\[0.3em]
\hline
Elektroauto (Antrieb) & $\square$ & $\square$ \\[0.3em]
\hline
Windkraftanlage & $\square$ & $\square$ \\[0.3em]
\hline
Faraday-Lampe (Schütteltaschenlampe) & $\square$ & $\square$ \\[0.3em]
\end{tabularx}

\end{merkkasten}

% ═══════════════════════════════════════════════════════════════
% ÜBUNG 2: ENERGIEKETTE
% ═══════════════════════════════════════════════════════════════

\begin{merkkasten}[Ü2: Energieumwandlung im Wasserkraftwerk]

\textbf{Ergänze die Energiekette:}

\vspace{3mm}

\begin{center}
\begin{tikzpicture}[node distance=1.2cm, every node/.style={font=\small}]
    \node[draw, fill=white, minimum width=3cm, minimum height=0.8cm] (lage) {\luecke{2.5cm}};
    \node[right=of lage, font=\Large] (pfeil1) {$\rightarrow$};
    \node[draw, fill=white, minimum width=3cm, minimum height=0.8cm, right=of pfeil1] (beweg) {\luecke{2.5cm}};
    \node[right=of beweg, font=\Large] (pfeil2) {$\rightarrow$};
    \node[draw, fill=white, minimum width=3cm, minimum height=0.8cm, right=of pfeil2] (elek) {\luecke{2.5cm}};

    \node[below=0.1cm of lage, font=\tiny] {(Wasser oben)};
    \node[below=0.1cm of beweg, font=\tiny] {(Turbine dreht)};
    \node[below=0.1cm of elek, font=\tiny] {(Strom fließt)};
\end{tikzpicture}
\end{center}

\vspace{1mm}

\begin{wortbank}
\centering
\textbf{Lageenergie} ~ | ~ \textbf{Bewegungsenergie} ~ | ~ \textbf{Elektrische Energie}
\end{wortbank}

\end{merkkasten}

% ═══════════════════════════════════════════════════════════════
% ÜBUNG 3: TRANSFER
% ═══════════════════════════════════════════════════════════════

\begin{merkkasten}[Ü3: Transfer -- Nachdenken]

\textbf{Beantworte in ganzen Sätzen:}

\vspace{2mm}

\textbf{a)} Ein Dynamo erzeugt bei normaler Fahrt eine Spannung von 0,5\,V. Was passiert ungefähr mit der Spannung, wenn du \textbf{doppelt so schnell} fährst?

\vspace{0.8cm}
\rule{\linewidth}{0.3pt}

\vspace{3mm}

\textbf{b)} Was passiert mit der Spannung, wenn du an der roten Ampel \textbf{stehen bleibst}? Erkläre warum!

\vspace{0.8cm}
\rule{\linewidth}{0.3pt}

\vspace{0.5cm}
\rule{\linewidth}{0.3pt}

\end{merkkasten}

% ═══════════════════════════════════════════════════════════════
% FUSSZEILE
% ═══════════════════════════════════════════════════════════════

\vfill

\begin{center}
\small\textcolor{mittelgrau}{Dieses AB begleitet den Lernpfad LP-09b. Nach Bearbeitung: Minitest MT-09b.}
\end{center}

\end{document}
